\documentclass{subfiles}
\begin{document}\label{sec:splay-trees}
    Introduced in \cite{SleatorT1985b} by \citeauthor{SleatorT1985b},
        splay trees are a variant of BSTs that are proved to be, in an amortized sense, optimal.
        In their paper the authors consider the following situation: 
        assume we are asked to perform a sequence of access operation on a set of items 
        selected from a totally ordered universe.
        The symmetry with the dictionary problem is obvious.

    The key observation of the authors is that, when accessing an item, 
        it is very likely that said item will be accessed soon,
        meaning that moving it toward the root will reduce the time required for the subsequent aceess operation.
        The idea, originally proposed in \cite{Bitner1979} by \citeauthor{Bitner1979},
        is improved in \cite{SleatorT1985b} to avoid the \(\ofO{n}\) time cost on
        arbitrary long sequences of operations. To achive so, they propose a new heuristic procedure 
        called splaying.

    Here we discuss splaying, how this is used to implement the dictionary operations, 
        and report some theoretical results regarding splay trees.
        See \cite{SleatorT1985b,Tarjan1983} for more details.

    Let \(x\) be the least recently accessed node, and let \(x.p\) denote \(x\)'s parent;
        we can divide the splaying procedure in the following three cases:
        \begin{description}
            \item [Case 1 (zig):  \(\mathbf{x.p}\) is the root.]
                Then, we simply left (resp. right) rotate on \(x\), depending if it was a right (resp. left)
                child of \(x.p\).

            \item [Case 2 (zig-zig): both \(\mathbf{x}\) and \(\mathbf{x.p}\) are left (resp. right) children.]
                We first right (resp. left) rotate on \(x.p\), making \(x.p\) a right child of \(x\),
                then we right rotate on \(x\).

            \item [Case 3 (zig-zag): \(\mathbf{x}\) is a right child and \(\mathbf{x.p}\) is a left child (or vice versa).]
                First, we left rotate on \(x\), making \(x.p\) be a left child of \(x\). Then,
                we right rotate on \(x\) and make the new parent a right child.
        \end{description}
        A visual representation of each case is shown in \Cref{fig:splay-example-splaying}.
        As for for RB trees, rotations are local operation on nodes, 
        therefore the order of the elements is left unchanged.

    The following analysis of the splaying procedure is based on the assumption that each to node 
        \(x\) is associated an arbitrary weight \(w(x)\). Additionally, we use the following notation:
        \(s(x)\) denotes the size of \(x\), i.e., the sum of the weight in the subtree rooted at \(x\),
        and by \(r(x)\) the rank of \(x\), defined as the logarithm of \(s(x)\).
        For a tree \(T\), we use \(r(T)\) to denote the rank of the root of \(T\).


    \begin{theorem*}[\cite[Lemma 1]{SleatorT1985b}]
        The amortized cost for splaying a tree \(T\) at a node \(x\) is at most \(3(r(T) - r(x)) + 1\).
    \end{theorem*}
    \begin{proof}
        If no rotations occur, the bound follow immediately.
            Hence, suppose at least a rotation is performed. Let \(s\) and \(s'\),
            \(r\) and \(r'\) be the weight and the rank before and after the rotation, respectively.
            We state that the amortized cost for the step is \(3(r'(x) - r(x)) + 1\) for Case 1,
            and \(3(r'(x) - r(x))\) for Case 2 and 3. Let \(y = x.p\) and \(z = y.p\), we have:
            \begin{description}
                \item [Case 1. ] One rotation occurs; the amortized time is therefore
                    \[\begin{aligned}
                        1 + & r'(y) - r(y) + r'(x) - r(x) && \text{since only \(y\) and \(x\) change rank} \\ 
                        \le & 1 + r'(x) - r(x) && \text{since \(r'(y) \le r(x)\)} \\ 
                        \le & 1 + 3(r'(x) - r(x)) && \text{since \(r'(x) \ge r(x)\)} \\ 
                    \end{aligned}\]

                \item [Case 2.] Two rotations occur; the amortized time is
                    \[\begin{aligned}
                        2 + & r'(z) - r(z) + r'(y) && \\ 
                        & - r(y) + r'(x) - r(x) && \text{since only \(z, y\) and \(x\) change rank} \\ 
                        \le & 2 + r'(z) + r'(y) - r(y) + r(x) && \text{since \(r(z) = r'(x)\)} \\
                        \le & 2 + r'(z) + r'(x) - 2r(x) && \text{since \(r'(x) \le r'(y)\) and \(r(y) \ge r(x)\)} 
                    \end{aligned}\]

                \item [Case 3. ] Two rotations are done, so the amortized time is 
                    \[\begin{aligned}
                        2 + & r'(z) - r(z) + r'(y) && \\ 
                            & - r(y) + r'(x) - r(x) && \\ 
                        \le & 2 + r'(z) + r'(y) - 2r(x) && \text{since \(r'(x) = r(z)\) and \(r(x) \le r(y)\)} 
                    \end{aligned}\]
            \end{description}

        The theorem follows by considering the maximum of the three cases. Hence,
            splaying at tree takes at most \(3(r(T) - r(x)) + 1\).
    \end{proof}
    \subfile{../TikZ/Figure 2 - Splay trees splaying}
\end{document}
